\section{The HCMO resource}
\textbf{Identity and size.} HCMO uses the persistent ontology IRI \nolinkurl{https://w3id.org/hcmo/ontology/hcm} and the stable term namespace \nolinkurl{https://w3id.org/hcmo/ontology/hcm#}. Domain-specific terms use four sub-namespaces beneath the same base: \nolinkurl{hcm/bio#}, \nolinkurl{hcm/env#}, \nolinkurl{hcm/obs#}, and \nolinkurl{hcm/tech#}. Published IRIs are not renamed when terminology changes; obsolete IRIs are retained in a compatibility module with deprecation and replacement metadata where a defensible replacement exists. Table~\ref{tab:resource-stats} summarises the release.

\begin{table}[t]
\centering
\scriptsize
\caption{Summary of the HCMO 0.3.0 release. \new{``Release'' counts all HCMO-namespace terms in the generated distribution, including deprecated compatibility terms; ``active'' excludes them. Reused external terms are counted separately.}}
\label{tab:resource-stats}
\begin{tabular}{@{}lrr@{}}
\toprule
Resource measure & Release & \new{Active} \\
\midrule
RDF triples in the canonical distribution & 1,397 & \\
HCMO classes & 60 & \new{33} \\
HCMO object properties & 68 & \new{32} \\
HCMO datatype properties & 75 & \new{37} \\
\new{External classes/properties reused for HCM semantics} & \new{25} & \\
\quad\new{SOSA 9, BFO 4, schema.org 4, SemTS 3, QUDT 3, IAO 1, PROV-O 1} & & \\
\new{External vocabularies reused for ontology metadata} & \new{6} & \\
Authored modules & 5 + compatibility & \\
Executable competency questions & 11 & \\
\bottomrule
\end{tabular}
\end{table}

\textbf{Modular organisation.} HCMO 0.3.0 is authored as five active domain modules. The minimal core centres on \nolinkurl{hcm:Enclosure}, \nolinkurl{hcm:MonitoredEnclosure}, enclosure dimensions, enrichment, and stable housing relations. The \emph{bio} module represents subjects, experimental groups, housing assignments, and study factors. The \emph{env} module represents environmental profiles, light cycles, gas profiles, measurement specifications, and environmental properties. The \emph{obs} module owns SOSA-style observations and their categorical, quantitative, behavioral, and tabular results. The \emph{tech} module represents sensors, actuators, hardware, software, and time-series resources. Two further modules complete the manifest: an upper-level presentation module supplying the BFO/IAO anchors, and a migration-only module that preserves deprecated HCMO 0.0.1 IRIs but is not a source of terms for new data. Figure~\ref{fig:domain-model} gives a compact view of the principal entities and relations.

\begin{figure*}[t]
\centering
\resizebox{\textwidth}{!}{%
\begin{tikzpicture}[
  entity/.style={draw, rounded corners=1.5pt, align=center, minimum height=9mm, minimum width=34mm, font=\scriptsize, inner xsep=4pt},
  rel/.style={-{Latex[length=1.8mm]}, font=\tiny},
  lbl/.style={fill=white, inner sep=1pt, font=\tiny},
  hlbl/.style={above, inner sep=1.5pt, font=\tiny},
  evidence/.style={draw, dashed, rounded corners=1.5pt, align=center, minimum width=34mm, font=\scriptsize, inner sep=4pt}
]
% Three rows on a fixed grid, so that no box or arrow crosses another:
% top = identity and observations, middle = housing and devices,
% bottom = time and environmental specifications.
\node[entity, fill=red!8]    (subject)     at (0,3)  {\texttt{hcm-bio:Subject}\\stable animal identity};
\node[entity, fill=blue!8]   (observation) at (11.2,3) {HCMO observation\\\texttt{sosa:Observation}};
\node[entity, fill=blue!8]   (result)     at (16.8,3)  {\texttt{hcm-obs:ObservationResult}\\\texttt{sosa:Result}};

\node[entity, fill=red!8]    (assignment) at (0,1.5) {\texttt{hcm-bio:HousingAssignment}\\time-bounded record};
\node[entity, fill=yellow!14] (cage)      at (5.6,1.5) {\texttt{hcm:MonitoredEnclosure}\\central physical hub};
\node[entity, fill=violet!9] (sensor)     at (11.2,1.5) {\texttt{hcm-tech:Sensor}\\\texttt{sosa:Sensor}};

\node[evidence]              (time)       at (0,0)  {\texttt{time:Interval}\\assignment validity};
\node[entity, fill=green!9]  (profile)    at (5.6,0)  {\texttt{hcm-env:EnvironmentProfile}};
\node[entity, fill=green!9]  (spec)       at (11.2,0) {\texttt{hcm-env:}\\\texttt{MeasurementSpecification}};
\node[entity, fill=green!9]  (quantity)   at (16.8,0) {\texttt{hcm-obs:QuantityValue}\\\texttt{qudt:QuantityValue}};

\draw[rel] (subject) -- node[lbl]{\texttt{hasHousingAssignment}} (assignment);
\draw[rel] (assignment) -- node[lbl]{\texttt{time:hasTime}} (time);
\draw[rel] (assignment) -- node[hlbl]{\texttt{assignedToEnclosure}} (cage);
\draw[rel] (cage) -- node[hlbl]{\texttt{monitoredBy}} (sensor);
\draw[rel] (sensor) -- node[lbl]{\texttt{sosa:madeObservation}} (observation);
\draw[rel] (observation) -- node[hlbl]{\texttt{sosa:hasFeatureOfInterest}} (subject);
\draw[rel] (observation) -- node[hlbl]{\texttt{sosa:hasResult}} (result);
\draw[rel] (cage) -- node[lbl]{\texttt{hasEnvironment}} (profile);
\draw[rel] (profile) -- node[hlbl]{\texttt{hasMeasurementSpec}} (spec);
\draw[rel] (spec) -- node[hlbl]{\texttt{hasSpecifiedValue}} (quantity);
\draw[rel, dashed] (quantity) -- node[lbl]{subclass} (result);
\end{tikzpicture}
}
\caption{The five-module HCMO model preserves animal identity, temporal housing,
environmental specifications, sensing events, and results as distinct entities.}
\label{fig:domain-model}
\end{figure*}


\textbf{Key modelling decisions.} This organisation preserves two boundaries that are important for HCM data. First, housing context is not reduced to the animal: a housing assignment is an explicit record linking a subject or group to an enclosure. Second, a physical sensor, the observation it performs, and the result it produces are separate entities. For example, \nolinkurl{hcm-tech:Sensor} is linked to a monitored enclosure, captures a \nolinkurl{sosa:ObservableProperty}, and may make a \nolinkurl{sosa:Observation}; the observation then identifies its feature of interest and result using SOSA relations. This supports provenance-sensitive queries without treating a device, event, and data value as the same thing. Environmental profiles, measurement specifications, and observations likewise use different relations: a profile composes specifications, each specification identifies the property and optional target quantity, and each observation uses the SOSA observed-property/result pattern. Housing membership is derived from time-bounded assignments at a stated time. Operational and calibration states are time-bounded records generated by explicit PROV activities rather than timeless booleans.

\textbf{Reused standards.} Physical subjects, enclosures, enrichments, sensors, actuators, hardware, and experimental groups are exposed under BFO material entity in the default end-user hierarchy. Records, specifications, profiles, software, and observation-result artifacts use the IAO information-content anchor. A separate optional developer profile restores BFO's source-faithful intermediate hierarchy and the more precise object-aggregate classification of experimental groups. Sensors, actuators, observations, results, and observed properties reuse SOSA classes and relations, and \nolinkurl{hcm-tech:captures} is explicitly a subproperty of \nolinkurl{sosa:observes}. HCMO normatively follows the 2017 SOSA/SSN Recommendation \cite{sosa}, pinned to an immutable copy of the 2017 ontology artifact; terms introduced by the developing later edition are not mixed into this release. Numeric dimensions, sampling rates, environmental targets, and observation results use the pinned QUDT 3.4.0 \nolinkurl{QuantityValue}, \nolinkurl{numericValue}, and \nolinkurl{hasUnit} pattern \cite{qudt}, with HCMO role properties stating what each quantity describes. HCMO also selectively reuses canonical SemTS 1.2.0 terms \cite{semts}: a location result table is a \nolinkurl{semts:TimeSeriesSegment} whose dimensions may be linked with \nolinkurl{semts:segmentDimension} to \nolinkurl{semts:DataDimension}; the earlier observation dimension restriction and \nolinkurl{semts:generated} relation were removed after semantic review. Schema.org supplies contributor and place exchange types, and the example ABox uses OWL-Time intervals. These are selective references to pinned artifacts, not claims that HCMO imports or reproduces the full external ontologies.

\textbf{Engineering and distribution.} HCMO is released as a \textbf{tool-consumable package}. A machine-readable manifest (\nolinkurl{hcmo.yaml}) declares the ontology's identity and version IRIs, the ordered source modules, the generated distributions, the SHACL shapes, the competency-question index, and the example datasets; its \emph{shape} is treated as an API and kept stable, so consumers discover paths from the manifest rather than hard-coding them. A build step merges the seven Turtle sources into a canonical, sorted serialisation that is byte-identical on re-run, published as Turtle, RDF/XML, and JSON-LD together with a flat term inventory (\nolinkurl{profile.json}) for synchronisation layers and interfaces, and a JSON-LD context (\nolinkurl{ontology/context.jsonld}) for application developers. \new{The presentation module adds a second \nolinkurl{owl:Ontology} node; metadata scanners should select the root header (version IRI, licence, DOI, creator ORCIDs) by its IRI.} A validation step parses every file, runs the SHACL shapes against positive and intentionally invalid examples, and executes every competency question against complete expected answers; it gates every pull request in CI, and a tag-triggered workflow publishes the release (Fig.~\ref{fig:pipeline}).

\begin{figure}[t]
\centering
\resizebox{\linewidth}{!}{%
\begin{tikzpicture}[
  node distance=4mm and 6mm,
  box/.style={draw, rounded corners=1.5pt, align=center, minimum height=8mm, font=\scriptsize, inner xsep=5pt},
  flow/.style={-{Latex[length=2mm]}, thick}
]
\node[box, fill=blue!8] (src) {Five source modules\\+ compatibility};
\node[box, fill=yellow!12, right=of src] (manifest) {Stable manifest\\\texttt{hcmo.yaml}};
\node[box, fill=green!10, right=of manifest] (build) {Deterministic build\\\texttt{build.py}};
\node[box, fill=green!10, right=of build] (dist) {TTL / OWL / JSON-LD\\profile inventory};
\node[box, fill=red!7, below=of build] (validate) {Parse + SHACL + 11 CQs\\audits + HermiT + ISA tests};
\node[box, fill=violet!8, right=of validate] (release) {CI gate\\tag + Zenodo + docs};
\draw[flow] (src) -- (manifest);
\draw[flow] (manifest) -- (build);
\draw[flow] (build) -- (dist);
\draw[flow] (dist.south) -- ++(0,-2.2mm) -| (validate.north);
\draw[flow] (validate) -- (release);
\draw[flow] (manifest) |- (validate);
\end{tikzpicture}
}
\caption{HCMO's manifest-driven, reproducible build and validation pipeline.}
\label{fig:pipeline}
\end{figure}


\textbf{Availability and sustainability.} The ontology and its terms dereference from the w3id IRI via HTTP 303 content negotiation to WIDOCO documentation \cite{widoco2017} (\url{https://dhuzard.github.io/HCMO/index-en.html}). Sources and distributions are openly available from the public GitHub repository and the Zenodo deposit\new{ (release \nolinkurl{v0.3.0}, DOI \nolinkurl{10.5281/zenodo.22208202})}. \reviewcomment
  {No hosted SPARQL service is part of HCMO 0.3.0.}
  {À confirmer avec Gaoussou et Konstantin.} The distributions can be loaded into any standards-compliant triple store. The resource is licensed CC BY 4.0, carries provenance on the root ontology header (creators with ORCIDs, version IRI), and provides a canonical citation (\nolinkurl{CITATION.cff}); a FAIR self-assessment is available in the repository. HCMO is maintained through public issues and pull requests; versioning follows SemVer mirrored by \nolinkurl{owl:versionIRI}, and terms are deprecated rather than deleted or re-minted to protect external references. \new{HCMO is registered in BioPortal (\url{https://bioportal.bioontology.org/ontologies/HCMO}).}\todo[color=orange!40]{DH: update the BioPortal entry with the final release before submission.}
